\documentclass[10pt,a4paper]{amsart}
\usepackage[margin=19mm]{geometry}
\usepackage{amsmath,amssymb,mathtools}
\usepackage[scr=rsfs,cal=euler]{mathalfa}
\usepackage{xfrac}
\usepackage[unicode,hidelinks,bookmarks=false]{hyperref}
\newcommand{\ie}{i.e.\ }
\newcommand{\cf}{cf.\ }
\newcommand{\resp}{respectively\ }
\newcommand{\Subsection}{\subsection}
\providecommand{\nobreakdash}{\nobreak\hbox{-}}
\setlength{\emergencystretch}{2em}
\multlinegap0pt
\usepackage{bm}
\usepackage{bbm}
\usepackage{dsfont}
\usepackage{xspace}
\usepackage{cancel}
\usepackage{mathtools}
\usepackage{amsthm}
\makeatletter
\@ifundefined{theorem}{\newtheorem{theorem}{Theorem}}{}
\@ifundefined{assumption}{\newtheorem{assumption}[theorem]{Assumption}}{}
\makeatother
\title[Spectral Analysis and Redistribution Thresholds for Cut-Cell]{Spectral Analysis and Redistribution Thresholds for Cut-Cell Finite-Volume Methods}
\author{Justo E. Karell}
\date{Source: arXiv:2607.28808v3 (CC BY 4.0)}
\begin{document}
\maketitle

\noindent\emph{Source and licence.} Spectral Analysis and Redistribution Thresholds for Cut-Cell Finite-Volume Methods. Version arXiv:2607.28808v3 (CC BY 4.0), distributed under the Creative Commons Attribution 4.0 International licence, as recorded in the arXiv licence metadata for this exact version and shown on the abstract page https://arxiv.org/abs/2607.28808v3. Full rights notice: licenses/arxiv-2607.28808-CC-BY-4.0.txt.\par\smallskip

\noindent\textit{Editorial scope.} Counted unit: the Theorem whose source label is \texttt{thm:proj} in the article (source0.tex lines 1152--1162, source label \texttt{thm:proj}), delivered together with the article's own complete original proof of it (source0.tex lines 1164--1270, one \texttt{proof} environment of 107 lines). No mathematics is written, repaired or re-derived: statement and proof are the article's own text line for line. Required context retained uncounted: ass:sb (lines 1035--1045); thm:singular (lines 1053--1074). Every definition, display and label of those lines is retained unchanged. Omitted from this package and not redistributed: 1--1034, 1046--1052, 1075--1151, 1163--1163, 1271--2776. Nothing inside a retained excerpt is omitted or altered: every retained body is delivered exactly as the article's lines are written, so no omission receipt of any kind accompanies this package. Citations: the retained excerpts carry no citation command at all, so no bibliography entry of the article is transcribed and no reference is redistributed. Reference adaptation: none was needed and none was made; every reference inside the delivered unit resolves inside it, so no unresolved reference or citation is produced. Printed slips kept literal: no printed word, symbol, hypothesis or number of the counted statement or of its proof was altered. Layout and class: the article's own class is \texttt{elsarticle} with packages including amsmath, amsthm, amssymb, mathtools, hyperref, graphicx, tikz; the wrapper is the lane's pinned \texttt{amsart} editorial wrapper at 10pt on A4, which restates the theorem-like environments the retained text needs and adds the article's own macro definitions that the retained text needs (none), with extra wrapper packages (bm, bbm, dsfont, xspace, cancel, mathtools). The delivered numbering is the unit's own. Assets: no retained span contains a figure environment, a graphics-inclusion command, a PDF-page inclusion, a table or a TikZ picture; no image file is copied into this package, the package declares no asset dependency and no image file is redistributed with it. Licence: version arXiv:2607.28808v3 of this article is distributed under the Creative Commons Attribution 4.0 International licence, as recorded in the arXiv licence metadata for this exact version and shown on the abstract page https://arxiv.org/abs/2607.28808v3. A full rights notice is delivered at licenses/arxiv-2607.28808-CC-BY-4.0.txt. Characters: every retained excerpt is pure ASCII, so the delivered engine, XeLaTeX with its default Latin Modern text font, reports no missing character.
\begin{assumption}[Asymptotic cut-cell block structure]
\label{ass:sb}
As $\alpha\to0$,
\[
    \|\alpha^pR_\alpha\|\to0,\qquad
    \|B_\alpha\|=O(\alpha^{-p}),\qquad
    \|C_\alpha\|=O(1),\qquad
    \|Q_\alpha\|=O(1),
\]
and $\Gamma$ is invertible.  The dimensions $m$ and $N$ are fixed.
\end{assumption}

\begin{theorem}[Spectrum generated by the leading cut-cell block]
\label{thm:singular}
Under Assumption~\ref{ass:sb}, $A_\alpha$ has exactly $m$ unbounded
eigenvalues, counted with algebraic multiplicity.  Let
$\gamma_1,\ldots,\gamma_m$ denote the eigenvalues of $\Gamma$, repeated
according to algebraic multiplicity.  The unbounded eigenvalues of $A_\alpha$
can be labeled $\mu_1(\alpha),\ldots,\mu_m(\alpha)$ so that
\begin{equation}
    \alpha^p\mu_k(\alpha)\longrightarrow-\gamma_k,
    \qquad k=1,\ldots,m.
    \label{eq:scaled-spectrum}
\end{equation}
All remaining eigenvalues are bounded.  If, in addition, $\Gamma$ is
diagonalizable and $R_\alpha=O(1)$, then the unbounded eigenvalues may be
matched with the eigenvalues $\gamma_k$ of $\Gamma$ so that
\begin{equation}
    \mu_k(\alpha)
    =
    -\frac{\gamma_k}{\alpha^p}+O(1).
    \label{eq:large-O1}
\end{equation}
\end{theorem}

\begin{theorem}[Convergence of the cut-cell invariant subspace]
\label{thm:proj}
Under Assumption~\ref{ass:sb}, let $\mathcal U_\alpha$ be the invariant subspace
associated with the $m$ unbounded eigenvalues.  Then
\begin{equation}
    \|P_{\mathcal U_\alpha}-P_{\mathcal C}\|=O(\alpha^p),
    \label{eq:proj}
\end{equation}
where $P_{\mathcal V}$ denotes the Euclidean orthogonal projector onto a
subspace $\mathcal V$.
\end{theorem}

\begin{proof}
Scale the matrix:
\[
 \widetilde A_\alpha:=\alpha^pA_\alpha
 =
 \begin{pmatrix}
     -\Gamma+\alpha^pR_\alpha & \alpha^pB_\alpha\\
     \alpha^pC_\alpha & \alpha^pQ_\alpha
 \end{pmatrix}.
\]
By Theorem~\ref{thm:singular}, the nonzero spectral cluster converges to
$-\sigma(\Gamma)$ and remains separated from the cluster at $0$.  Choose a
fixed positively oriented contour $\mathfrak C$ enclosing
$-\sigma(\Gamma)$ and excluding $0$.  The Riesz projector
\[
    \Pi_\alpha
    =
    \frac{1}{2\pi i}
    \int_{\mathfrak C}
    (zI-\widetilde A_\alpha)^{-1}\,dz
\]
has range $\mathcal U_\alpha$ and rank $m$.

Write
\[
 zI-\widetilde A_\alpha
 =
 \begin{pmatrix}
 X_\alpha(z) & -\alpha^pB_\alpha\\
 -\alpha^pC_\alpha & Y_\alpha(z)
 \end{pmatrix},
\]
where
\[
 X_\alpha(z)=zI+\Gamma-\alpha^pR_\alpha,
 \qquad
 Y_\alpha(z)=zI-\alpha^pQ_\alpha.
\]
Since $\mathfrak C$ stays a positive distance from $0$,
$Y_\alpha(z)^{-1}=O(1)$ uniformly on $\mathfrak C$.  The Schur complement
\[
 \mathcal S_\alpha(z)
 =
 X_\alpha(z)
 -
 \alpha^{2p}B_\alpha Y_\alpha(z)^{-1}C_\alpha
\]
satisfies
\[
 \mathcal S_\alpha(z)
 =
 zI+\Gamma+o(1)
\]
uniformly on $\mathfrak C$, because
$\alpha^{2p}B_\alpha Y_\alpha(z)^{-1}C_\alpha=O(\alpha^p)$.
Hence $\mathcal S_\alpha(z)^{-1}=O(1)$ uniformly on the contour.

The block inverse formula now gives
\[
 (zI-\widetilde A_\alpha)^{-1}_{21}
 =
 Y_\alpha(z)^{-1}\alpha^pC_\alpha\mathcal S_\alpha(z)^{-1}
 =
 O(\alpha^p),
\]
and
\[
 (zI-\widetilde A_\alpha)^{-1}_{22}
 =
 Y_\alpha(z)^{-1}
 +
 Y_\alpha(z)^{-1}\alpha^pC_\alpha\mathcal S_\alpha(z)^{-1}
 \alpha^pB_\alpha Y_\alpha(z)^{-1}.
\]
The second term is $O(\alpha^p)$.  The first term,
$Y_\alpha(z)^{-1}$, is analytic inside $\mathfrak C$ because the spectrum of
$\alpha^pQ_\alpha$ remains near $0$, outside the contour.  Its contour
integral is therefore zero.  It follows that
\[
    (\Pi_\alpha)_{21}=O(\alpha^p),
    \qquad
    (\Pi_\alpha)_{22}=O(\alpha^p).
\]
Similarly,
\[
    (\Pi_\alpha)_{11}=I_m+o(1),
\]
because the contour encloses all $m$ eigenvalues of $-\Gamma$.

For small $\alpha$, $(\Pi_\alpha)_{11}$ is invertible.  The $m$ vectors
$\Pi_\alpha(x,0)^\top$ therefore span $\operatorname{range}(\Pi_\alpha)$ and
show that
\[
    \mathcal U_\alpha
    =
    \{(x,K_\alpha x):x\in\mathbb C^m\},
    \qquad
    K_\alpha
    =
    (\Pi_\alpha)_{21}(\Pi_\alpha)_{11}^{-1}
    =
    O(\alpha^p).
\]
The standard orthogonal-projector formula for the graph of $K_\alpha$ then
gives
$\|P_{\mathcal U_\alpha}-P_{\mathcal C}\|=O(\|K_\alpha\|)=O(\alpha^p)$.
\end{proof}

\end{document}
